% ═══════════════════════════════════════════════════════════════════════════
\chapter{Introduction}
\label{ch:introduction}

\noindent\textit{(by S.J.\ Pogue, P.\ Mantle, A.\ McPherson and R.\ Kröbel)}

% ───────────────────────────────────────────────────────────────────────────
\section{Purpose of the Holos model}
\label{sec:purpose}

The Holos model is a software application developed and provided by Agriculture
and Agri-Food Canada (AAFC) to estimate greenhouse gas (GHG) emissions and soil
C changes from Canadian farming systems. The development of the Holos model
drives innovation and ingenuity for the benefit of all Canadians through the
adoption of scientific research that develops new knowledge and technologies,
transferring them to the agriculture and agri-food sector in the form of a
software application. The model's development focuses on enhancing the
environmental sustainability of Canadian farming systems (while attempting to
capture market opportunities) by testing, validating, and providing
sustainability attributes for Canadian agricultural production systems.

The development of the Holos model aligns with AAFC's agro-ecosystem resilience
strategy, which aims to reduce the environmental footprint of agricultural
production, and thus also interlinks closely with other AAFC sector strategies.
For this purpose, the model combines country-specific emission factors (EFs) and
algorithms in a whole-farm approach that encompasses all components of a farm
system, which can include different common livestock types such as beef, dairy,
pork, and poultry, but also less common types such as sheep, bison, horses,
etc., annual crops (e.g., wheat, canola, corn, various legumes and other small
grain crops), pastures and grasslands. The model can also provide estimates of
fuel and electricity demand (e.g., irrigation pumps) from farming operations
and infrastructure.

The model estimates direct nitrous oxide (N\textsubscript{2}O), methane
(CH\textsubscript{4}), and carbon dioxide (CO\textsubscript{2}) losses from
the different farm components, and indirect losses as a consequence of ammonia
(NH\textsubscript{3}) volatilization and nitrate
(NO\textsubscript{3}\textsuperscript{--}) leaching. In addition to these, the
model includes `upstream' estimates for emissions created through the production
of farm inputs (e.g., CO\textsubscript{2} losses from fertilizer, pesticide,
and electricity production). This allows for the estimation of GHG emissions
from farming systems, as well as commodity-based `cradle to farm-gate'
assessments that can be part of life cycle analyses.

While there is the option to test the results of scientific experiments in a
whole-farm context (under `what is?' scenarios), it is the aim of the model to
allow for the exploration of alternative management scenarios that may aid the
decision-making processes of producers and policy makers alike (answering `what
if?' questions). Holos V4 aims, thus, to provide an easy-to-use model for
multiple users (including producers, researchers and policy makers), allowing
them to investigate the impacts of on-farm practices, management decisions and
policy interventions on GHG emissions.

With the development of version 4 (V4) of the model, the Holos team recognizes
the need to further facilitate the holistic assessment of farms with additional
components (e.g., on-farm wetlands, riparian buffers, woodlots, and renewable
energy generation), as well as the need to provide additional capabilities
(e.g., economics and ecosystem services modules). Although these components are
not featured in Holos V4, the future collaborative development of Holos (V5 and
thereafter) has been facilitated by the open source availability of the V4
model code. Thus, we are continually searching for opportunities to involve
the expertise of other groups in- and outside of AAFC to facilitate the
improvement and expansion of future versions.

% ───────────────────────────────────────────────────────────────────────────
\section{Target audience}
\label{sec:audience}

There is not one single target user group for Holos, as the model aims to serve
the needs of users from a variety of backgrounds and with different objectives.
These objectives may be scientific in nature, related to policy making, or the
user may wish to explore the impacts of alternative farm management decisions.
In the latter sense, the model can also be used as an educational tool, opening
up the user community to also include educators. The development team aims to
increase the capabilities and usability of the model by providing options for
collaborative development for computer science experts or multi-media artists.

Scientific use of the model concentrates primarily on the (re)assessment of
results from scientific experiments, comparing emission measurements to the
model outputs, or testing the effects of these experiments in the context of
overall farm emissions. This has led to a number of peer-reviewed publications
that have examined the emission sources in beef and dairy production systems,
comparing past and present performance, and investigated the effect of feed
ingredient choice on overall farm emissions. The latter involves the estimation
of crop production emissions and carbon (C) storage potentials from crop
production choices (e.g., reduction of tillage or planting of perennials). For
scientific use, we expect detailed data inputs that reflect the experimental
setup.

Policy makers are likely to utilize the model for policy implementation, where
farmers (landowners) input their farm operation data into the model and use the
outputs (potentially through a certified third party) as confirmation that the
policy has been fulfilled. This requires certainty that the model algorithms
work reliably (i.e., by alignment with national metrics), that the specific
policy option can be modelled, and that the farmer (landowner) has all required
input data available. Another option for the policy maker is, however, to test
policies and their effects before implementation. For that purpose, the model
requires appropriate regionally representative default farming systems (with all
of the necessary input parameters) to test the policy under different settings.

For farmers, the model could be useful to self-explore where emissions come from
on their farm, and how they compare in magnitude. In addition, they can explore
management practices that lower emissions from the farm, altering individual
practices or the entire farm system. Productivity changes due to management
decisions are not included in the model, as the estimation of crop or animal
growth would require additional inputs that are not readily available, and would
also lead to considerable uncertainty (note that a simple economic model for
cropping operations attempts to provide cost estimates for changing management,
but is limited in scope and operability). Regardless of these limitations,
farmers have used the model for self-assessments (`what if?') for marketing
purposes.

The model has further application in agricultural education, particularly in
courses concerned with sustainability. This serves both to showcase to students
how interconnected production processes are in farming systems, and also to
identify GHG sources in farming and the best options to reduce them effectively.
The software development in itself has also proven to be an excellent teaching
opportunity for individual computer science students who want to experience
industry standard programming and code maintenance. Last but not least, the team
is looking to partner with multimedia and arts students to improve the design
and form of model inputs and outputs.

% ───────────────────────────────────────────────────────────────────────────
\section{Principles of model development}
\label{sec:principles}

Development of the Holos model takes place in two phases: the development and
combination of algorithms, and the transferral of these algorithms into software
code in conjunction with the design of the interface. This is not necessarily a
one-way street, as both interface design and software development may trigger
additional development (mainly due to opportunity rather than necessity). In any
case, the Holos development team follows the two principles of transparency and
scientific reliability in order to foster open and collaborative model and
software development.

Scientific reliability is at the core of Holos algorithm development and means
that the team attempts to build the model entirely out of algorithms and
relationships that have been published in peer-reviewed articles, which have,
where possible and/or required, a Canadian context. This approach does mean that
the model will not be abreast of current scientific developments in the field,
but rather that new knowledge will only be adopted when scientific consensus is
built. Scientific reliability is also achieved through alignment with other
metrics and assessment methods, with the National GHG inventory and AAFC's
Agri-environmental indicators being important benchmarks and algorithm
providers. Lifecycle analysis standards (ISO 14044) also provide important
guidance on how model results should be summarized. Scientific reliability and
the opportunity to trace each algorithm to its published origin is also the
foundation of the transparency principle.

In an attempt to avoid Holos model outputs being perceived as the result of
`black box' simulations, the Holos development team releases algorithm documents
that detail each and every algorithm incorporated into the model (equation
numbers in the document will match equation numbers in the software code) as
well as the lookup values and default data used to populate the model. Algorithm
documents track the progress of consecutive model versions, and in principle
provide the basis for re-building the model in any other (software) environment.
To further expand the transparency and adaptability of the model, the model's
software code is shared in an open source environment on the
\href{https://github.com/holos-aafc/Holos}{Holos GitHub repository}. This is
with the goal of sharing the algorithm interconnectedness that is hard to relay
in written documents, but also to open the model's software code for
collaborative additions by other groups, and for adaptation of the model by
other groups for their specific objectives and requirements.

The Holos development team pledges to create and share a framework for
collaborative development with respect to software architecture. However, the
determination of what is included in subsequent Holos model versions remains
with the Holos development team for the foreseeable future.

% ───────────────────────────────────────────────────────────────────────────
\section{Workshops and outreach}
\label{sec:workshops}

Modelling agricultural systems is a way to test our scientific understanding of
processes taking place in these systems, to provide context to the results of
scientific experiments, to understand and teach the interconnectedness of farm
systems, to explore options and capabilities to influence the system and its
outputs, and to guide and inform future research. Despite these benefits,
modelling these systems doesn't come naturally to any of our potential user
groups, and the team, therefore, provides multiple options to learn more about
modelling and the model. These options range from organizing an annual
workshop/conference (including a model training session), providing model
training as invited course or seminar speakers, providing training upon request
or online, developing and distributing training materials, hosting trainees for
model training or development purposes, as well as conducting and publishing our
own modelling studies.

Our annual event has evolved from an initial afternoon Holos release workshop
with 10 participants to a three-day `Sustainability of Canadian Agriculture'
conference with several hundred attendees, where we bring together scientists,
academics, farmers, farmer organizations, and policy makers to debate how
Canadian farming can continue into the future, and, in this sense, how the
Holos model and its team can help and support the community. During this
conference, our half-day Holos training session has attracted over a hundred
participants in recent years.

Model training has been provided on demand on multiple occasions, mostly upon
invite from university partners, although the provincial government has also
requested model training. Typically, the Holos team appreciates a formal invite
($>$2 months before the event) to organize travel (contact
\href{mailto:holos@agr.gc.ca}{holos@agr.gc.ca}). Training can be held for up
to 50 participants (it is advisable that trainees work in pairs), and computers
(Windows platform) should be available with the Holos model pre-installed. It
is helpful if trainers other than the Holos team are familiar with the training
to assist participants. Training can also be provided through an online platform,
although troubleshooting is more limited and there is less opportunity for
feedback.

One element that can improve online training, in particular self-training, are
our Holos training materials, which include step-by-step pdf training guides and
a series of video tutorials (available on our
\href{https://github.com/holos-aafc/Holos/tree/main}{GitHub page}). The team is
available to host trainees at the Lethbridge Research and Development Centre and
also offers opportunities for students as part of the Canadian universities
Co-op student program.

Last but not least, the Holos team conducts its own model application studies
that aim to demonstrate how the model can be utilized. Publications relating to
model development and model application are listed on our
\href{http://www.agr.gc.ca/eng/scientific-collaboration-and-research-in-agriculture/agricultural-research-results/holos-software-program/?id=1349181297838}{Holos download page},
but we have limited resources to find publications by other groups who utilized
the model. We would therefore appreciate if you could send us a copy!

% ───────────────────────────────────────────────────────────────────────────
\section{Structure of this document}
\label{sec:structure}

The Holos algorithm document provides a detailed description of the algorithms
incorporated into the Holos model. It begins with the land management section
(\textbf{Chapter~2}, also the recommended first step in model input for users),
detailing at first the calculations for carbon change using the Introductory
Carbon Balance Model (ICBM), and subsequently the IPCC Tier~2 carbon model
(IPCC 2019), the current National GHG inventory approach and, thus, the default
option in Holos. The final carbon section is the updated shelterbelt component,
from where we plan to continue into the water budget sub-model (this will be
added after the release of V4). With the nitrogen (N) model, we describe the
multi-year N budget approach (compliant with the National GHG inventory as far
as the emissions estimates are concerned).

In the livestock sections (\textbf{Chapters~3, 4 and 5}), we separate beef
cattle, dairy cattle, swine, sheep, poultry and other livestock, using an IPCC
Tier~2 methodology for beef and dairy cattle and some poultry calculations, and
a mostly IPCC Tier~1 methodology for all other livestock types.

\textbf{Note:} throughout the document the term ``livestock type'' refers to
the broad livestock categories included in the model, e.g., beef cattle, dairy
cattle, swine, etc., while the term ``animal group'' refers to individual
sub-groups of animals within each livestock type, e.g., beef cows, dairy
heifers, sows, rams, turkey hens, etc.

In these sections, the model estimates emissions from livestock and the manure
they produce, tracking C and N flows from excretion through to land application.
\textbf{Chapter~4} also includes an anaerobic digestion (AD) component, into
which the user can `load' manure and crop residues, with this sub-model tracking
C and N losses throughout the AD process and the production of biogas, heat and
electricity. In \textbf{Chapter~6} the energy requirement calculations (and
emissions) are detailed. The economics section is still under development, as we
determine how to move forward from the initial economics component in Holos V3
(also to be added after the release of V4). The end of the document details the
summation of the results, uncertainty estimates and reporting considerations,
before finishing with the appendices.

% ───────────────────────────────────────────────────────────────────────────
\section{Required input parameters}
\label{sec:inputs}

\noindent\textbf{General considerations:}

\begin{itemize}[leftmargin=1.5em]
    \item \textbf{Metric or imperial units of measurement are chosen at the
          start, and apply to the complete run.} \textit{Note: the imperial
          option is currently unavailable but will be re-enabled in a future
          version of Holos.}

    \item \textbf{There is a French and an English version of the interface.}

    \item \textbf{All information entered saves instantaneously and
          continuously.}

    \item \textbf{The user can select each farm component multiple times.}

    \item \textbf{The default soils data are limited to the areas designated
          as containing agricultural land by Statistics Canada.}

    \item \textbf{When adjusting inputs in an already saved farm, press
          `reload data from previous screen' on the Details screen to update
          the results.}

    \item \textbf{In the options, model defaults can be reset to the latest
          lookup-table values in the settings menu.}
\end{itemize}

% ·····················································
\subsection{Map location}
\label{subsec:map}

The first required model input is the selection of the farm's location. The
user first selects the relevant Canadian province, and then the specific
longitude and latitude of the farm by clicking on the relevant location in the
displayed map. Selecting the specific farm location also allows Holos (as a
default) to download a location-specific historic daily weather dataset from
\url{https://power.larc.nasa.gov/data-access-viewer/}. These climate data are
required for the ICBM (detailed in \textbf{Section~2.1}), the IPCC Tier~2
carbon model (\textbf{Section~2.2}), the calculation of N\textsubscript{2}O
emissions from agricultural soils (\textbf{Sections~2.5, 2.6, and 2.7}), and
the estimation of emissions from livestock (\textbf{Chapters~3, 4 and 5}).

Each farm location will fall within a specific Soil Landscapes of Canada (SLC)
polygon. The choice of SLC polygon allows the lookup of commonly present soil
types in the region
(\url{https://open.canada.ca/data/en/dataset/5ad5e20c-f2bb-497d-a2a2-440eec6e10cd}),
permitting the user to choose the best fitting soil type (where more than one
exists within a single SLC polygon), which is then applied as a default
throughout the farm.

\textbf{Note:} the user can specify soil data for individual fields on the
`Soil' tab for the relevant field -- this will override the default soil data
used for the entire farm.

Also in the Map Location screen, on a different tab, the appropriate plant
hardiness zone (PHZ) can be selected
(\url{http://www.planthardiness.gc.ca/?m=1}), if more than one is present in
the selected polygon. The PHZ is used for specific shelterbelt soil C and root
growth lookup tables, with the exception of the province of Saskatchewan (SK),
where lookup tables are present for each ecodistrict (this is one case where
the province selection triggers a province-specific lookup).

% ·····················································
\subsection{Climate data}
\label{subsec:climate}

The resolution of the climate data required by Holos varies depending on the
model component. Daily climate data are needed for the calculation of
NH\textsubscript{3} losses and animal performance, as well as in the land
management component to calculate the soil moisture, but the model is set up to
run with monthly climate normals if daily data are not available.

For the ICBM, daily temperature and precipitation data are required for the
estimation of soil temperature and as input to the soil moisture model;
additionally, daily potential evapotranspiration data are required. Research
conducted by Martel et al.\ (2018) found that the Turc (1961) equation was best
suited to estimate actual evapotranspiration from Canadian farming systems,
compared with measured data from an Eddy-flux system. These equations (in
response to the relative humidity) are given below.

\bigskip
\noindent\textbf{For $RH \geq 50\%$:}

\begin{equation}
    ET_r = 0.013 \times \frac{T_{\mathrm{mean}}}{T_{\mathrm{mean}} + 15}
           \times \bigl(23.8856 \times R_s + 50\bigr)
    \tag{Eq.\ 1.6.2.1}
\end{equation}

\noindent\textbf{For $RH < 50\%$:}

\begin{equation}
    ET_r = 0.013 \times \frac{T_{\mathrm{mean}}}{T_{\mathrm{mean}} + 15}
           \times \bigl(23.8856 \times R_s + 50\bigr)
           \times \left(1 + \frac{50 - RH_{\mathrm{mean}}}{70}\right)
    \tag{Eq.\ 1.6.2.2}
\end{equation}

\noindent\textit{Turc (1961) after Alexandris et al.\ (2008)}

\begin{multline}
    K_c = a + b\!\left(\frac{T_{\max}-T_{\min}}{2} - T_{\mathrm{base}}\right)
            + c\!\left(\frac{T_{\max}-T_{\min}}{2} - T_{\mathrm{base}}\right)^{\!2} \\
            + d\!\left(\frac{T_{\max}-T_{\min}}{2} - T_{\mathrm{base}}\right)^{\!3}
            + e\!\left(\frac{T_{\max}-T_{\min}}{2} - T_{\mathrm{base}}\right)^{\!4}
    \tag{Eq.\ 1.6.2.3}
\end{multline}

\noindent For grasses, a non-linear function to calculate $K_c$ is used:

\begin{equation}
    K_c^{(\mathrm{grass})} = a \times
        \left(1 - e^{-b\left(\tfrac{T_{\max}-T_{\min}}{2} - T_{\mathrm{base}}\right)}\right)
    \tag{Eq.\ 1.6.2.4}
\end{equation}

\begin{equation}
    ET_c = K_c \times ET_r
    \tag{Eq.\ 1.6.2.5}
\end{equation}

\noindent where

\begin{center}
\begin{tabular}{@{}ll@{}}
\toprule
\textbf{Symbol} & \textbf{Definition} \\
\midrule
$ET_r$              & Reference evapotranspiration (mm day\textsuperscript{-1}) \\
$ET_c$              & Crop-specific evapotranspiration (mm day\textsuperscript{-1}), \\
                    & without consideration of soil water availability \\
$T_{\mathrm{mean}}$ & Mean daily temperature (\textdegree C) \\
$R_s$               & Solar radiation (MJ m\textsuperscript{-2} day\textsuperscript{-1}) \\
$RH_{\mathrm{mean}}$& Relative humidity (\%) \\
$K_c$               & Growing degree day driven crop coefficients \\
$T_{\mathrm{base}}$ & 5 \textdegree C (Martel et al.\ 2018) \\
$a,b,c,d,e$         & Coefficients required to calculate $K_c$ (Table~1) \\
\bottomrule
\end{tabular}
\end{center}

\bigskip
In order to reduce data input requirements for the user, the Holos model links
to a repository of daily climate data from the NASA POWER data access viewer:
\url{https://power.larc.nasa.gov/data-access-viewer/}.

For future scenarios, Holos calculates monthly climate normals for the
projection period based on daily NASA data (when there is an internet
connection) or using in-built Holos climate normals for each SLC polygon (when
no internet connection is available). Using the NASA climate data, Holos
calculates climate normals for the 2000--current year period as a default;
using the SLC-specific climate normals, the calculation period is 1990--2017.
The climate normals are calculated as follows:

\begin{equation}
    \mathrm{Climate\_Normal}_{(\mathit{type},\,\mathit{period},\,\mathit{month})}
    = \frac{\displaystyle\sum_{i=\mathit{Start}_{\mathit{year}}}^{\mathit{End}_{\mathit{year}}}
            \sum_{i=\mathit{Start}_{\mathit{month}}}^{\mathit{End}_{\mathit{month}}}
            \mathrm{Climate\_value}_{(\mathit{type},\,\mathit{day},\,\mathit{month},\,\mathit{year})}}
           {(\mathit{End}_{\mathit{month}} - \mathit{Start}_{\mathit{month}}) \times 30}
    \tag{Eq.\ 1.6.2.6}
\end{equation}

\noindent where

\begin{center}
\begin{tabular}{@{}lp{9cm}@{}}
\toprule
\textbf{Symbol} & \textbf{Definition} \\
\midrule
$\mathrm{Climate\_Normal}_{(type, period, month)}$
    & 30-year average of daily climate values for a certain \textit{type},
      \textit{period}, and \textit{month} \\
\textit{Period}
    & Range of years for which climate normals are calculated (Table~2) \\
\textit{Month}
    & Range of months for which climate normals are calculated (Table~2) \\
\textit{Type}
    & Maximum temperature (\textdegree C), minimum temperature (\textdegree C),
      mean temperature (\textdegree C), precipitation (mm), crop
      evapotranspiration (mm), relative humidity (\%), solar radiation
      (MJ m\textsuperscript{-2}) \\
$\mathit{Start}_{\mathit{year}}$
    & Year in which a climate normal period starts (Table~2) \\
$\mathit{End}_{\mathit{year}}$
    & Year in which a climate normal period ends (Table~2) \\
$\mathit{Start}_{\mathit{month}}$
    & Julian day on which a month starts (Table~2) \\
$\mathit{End}_{\mathit{month}}$
    & Julian day on which a month ends (Table~2) \\
\bottomrule
\end{tabular}
\end{center}

% ·····················································
\subsection{Land management}
\label{subsec:land}

\subsubsection{Field inputs}

\begin{itemize}[leftmargin=1.5em]
    \item Crop type grown (incl.\ cover crop) and field size are essential
          inputs, as these are needed to estimate N inputs from aboveground
          (i.e., straw and product (harvest losses)) and belowground (i.e.,
          roots and root exudates/extra-root material) biomass to the soil.
          To carry out the necessary calculations, the model uses the input
          data together with in-built default data. If the user wants to
          specify a (repeating) cropping sequence/crop rotation, the crops
          planted in previous years also need to be specified.

    \item Crop yield, amount of fertilizer applied, and residue return should
          be provided by the user if possible. In the absence of user-defined
          inputs for these parameters, the model will use default data,
          although these are only approximate representations:
          \begin{itemize}
              \item Default crop yields are based on yields from the Small
                    Area Data (SAD) database for the specified crop and year
                    (\url{https://www150.statcan.gc.ca/t1/tbl1/en/tv.action?pid=3210000201}).
                    When the user selects `Average Yield' as the Yield
                    Assignment Method, Holos calculates this as the average
                    of the yield for all crops in the rotation as specified
                    under Step~3 in the Component Selection screen.
              \item Default fertilization rates (N) are calculated from crop
                    N demand (based on the crop yield), assuming a
                    fertilization efficiency of 50\% (thus by default
                    doubling the crop N demand).
              \item Manure N and C additions are estimated from the livestock
                    components or from literature derived averages.
              \item The residue return default for roots is 100\% for all
                    annual crops (except root crops), and straw return is
                    assumed to be 100\% (all straw remaining in the field).
                    For perennial crops, 30\% of the root biomass is assumed
                    returned to the soil in all years except the final year
                    of the stand, when this value is 100\%. Harvest losses
                    are assumed to be 2\% for annual crops and 35\% for
                    perennial crops.
          \end{itemize}

    \item Tillage regime, irrigation amount, and pesticide application rates
          are minor inputs for GHG emission estimates, as their impact is
          primarily driven by the emissions caused through fuel use. However,
          they are major cost elements in a farm system and will have more
          importance in the economics component. Atmospheric N deposition,
          N fixation, soil test N and fertilizer efficiency (to reduce the
          estimated N fertilization rate), and the moisture content of the
          crop at harvest can also be adjusted by the user.

    \item Crop-specific biomass fractions and N concentrations are `hidden'
          inputs that can be overwritten if specific or better data are
          available, but these are considered expert level changes and making
          these adjustments is not recommended for the average model user.
\end{itemize}

\subsubsection{Crop rotation}

The crop rotation component has the same input requirements as the field
component, with the difference being that each added rotation phase (i.e.,
year) will automatically add one field within the same rotation, but with the
crop sequence shifted by one year. This is done to maintain the principle that
each rotation phase is present each year. The crop management practices
specified are automatically applied to all fields for the entire duration of the
simulation, although these can be modified later in the `Timeline' screen.

\subsubsection{Shelterbelt}

For the shelterbelt component, the current state of C accumulation can be
calculated using row length, number of trees and their average circumference,
as long as the shelterbelt species is specified. The age of the shelterbelt can
also be defined. The model will provide estimates of the amount of C stored
within the tree and C stored in the soil below. Projections into the future
presume non-limited (ideal) growth if planted new, but if a below-average
circumference is specified for an existing shelterbelt, that reduction will
perpetuate into future C accumulation estimates.

% ·····················································
\subsection{Livestock management}
\label{subsec:livestock}

The livestock management interface for all animal groups follows the same
principle. In the first step, the user defines how many animal groups are in
their system. Animals can be partitioned into groups with similar management
practices, but if so desired, the management for individual animals can be
specified. In the second step, users define the management periods for their
system (by animal group), again with the option to detail management practices
for individual days if desired. The third step permits the user to specify the
number of animals, their start weight and average daily gain. Model default
values for the housing and manure management systems can be overwritten.
Furthermore, diet quality can be defined, either by using the provided defaults,
or by designing a new diet using the Custom Diet Creator. When choosing pasture
as the housing system, an additional dropdown will enable the user to `place'
the animals onto one of the fields that have been added to the farm. This will
trigger the model to calculate excreted manure as a C and N addition to that
field.

\subsubsection{Beef}

The beef production system provides three components: the cow-calf component
that is pre-setup to cover a full year, and the (at least initially)
shorter-term backgrounding and finishing components. The cow-calf component has
presets for cow/calf groups, replacement heifers and bulls, and preset
management periods for winter feeding, summer grazing and extended fall grazing.

Preset animal numbers, weights and weight gain, housing system, manure
management and diets are also present for all beef components, but note that
only one management period is entered as default for the backgrounding and
finishing components. If the user wishes to explore emissions over a full year,
they will have to add management periods to represent subsequent cattle
generations.

The user should override these defaults as needed, with the primary focus
(regarding the calculation accuracy of GHG emissions) on animal numbers and
animal weights (and gain), followed by diet specifications, and the housing and
manure management options. Selecting ``Show Additional Information'' in the
``Housing'' tab also grants access to the bedding material calculator, and in
the ``Manure'' tab this grants access to several EFs used in the model
calculations. The latter should not be changed without scientific evidence.

\subsubsection{Dairy}

The dairy production system does not feature different components, but instead
lists the different animal groups within a single component. However,
group-specific management periods and practices are provided and can (and
should) be overwritten by the user where locally-specific data are available.
In contrast to the beef components, milk production, milk fat content and milk
protein should be specified by the user as this information will impact the
energy requirement calculation, and thus the enteric CH\textsubscript{4}
estimates.

\subsubsection{Pork}

The pork/swine production system offers different components (or combinations
thereof) that represent typical Canadian farming systems (Farrow-to-Wean,
Iso-Wean, Grower-to-Finish, and Farrow-to-Finish). The user should select the
components that best represent their system. Only the Farrow-to-Wean and
Farrow-to-Finish include Sow and Boar groups as default, but animal groups can
be added to any component as needed. As it is common to switch diets based on
weight during the different stages of growth, the default management periods are
set up to represent expected changes in diet as animals age.

Note that the enteric CH\textsubscript{4} emissions from pork are calculated
using a per animal emission rate, specific to different animal groups and weight
categories. These EFs can be adjusted when the user selects the ``Show
Additional Information'' option, although these should only be changed based on
scientific evidence. By default, the model assumes that all animals are housed
in barns but there is also the option to place the animals on pasture. Diet
formulations can also be changed, but as mentioned above, this will not impact
the CH\textsubscript{4} emission rate until a method is developed that can
account for such changes. However, changes in diet do impact C and N excretion
in manure and subsequent emissions estimates.

\subsubsection{Poultry}

The poultry components represent different operations commonly present in Canada
(Pullet Farm, Chicken Multiplier Breeder, Chicken Multiplier Hatchery, Chicken
Meat Production, Chicken Egg Production, Turkey Multiplier Breeder, and Turkey
Meat Production), but these are more limited than the other livestock components
due to the lack of a model that can take diet information into account. The
poultry model therefore does not require feed information at this time. Animal
numbers, housing, and manure management are inputs that can be adjusted.

\subsubsection{Sheep}

The sheep components include Rams, Lambs and Ewes, and a Sheep Feedlot. Input
requirements are similar to the beef cattle and dairy components. However, due
to a lack of available data, Holos does not currently include Canada-specific
default management periods for these animals. For the diets, the user can choose
the preset default diets or define custom diets in the Custom Diet Creator.

\subsubsection{Other livestock}

All other livestock components use the IPCC Tier~1 (IPCC 2006, 2019)
methodology (default emission rates per animal), which is why only animal
numbers and manure management information are needed. Manure excretion rates are
estimated, so animals can be placed on pasture, but currently no biomass
consumption rates are available for these animal groups.

% ───────────────────────────────────────────────────────────────────────────
\section{References}
\label{sec:references}

\begin{description}[leftmargin=0pt, style=nextline]

    \item[IPCC, 2006]
        \textit{2006 IPCC Guidelines for National Greenhouse Gas Inventories.}
        Edited by S.\ Eggelston, L.\ Buendia, K.\ Miwa, T.\ Ngara, K.\ Tanabe.
        Published by the Institute for Global Environmental Strategies (IGES)
        for the IPCC. ISBN 4-88788-032-4.\\
        \url{https://www.ipcc.ch/report/2006-ipcc-guidelines-for-national-greenhouse-gas-inventories/}

    \item[IPCC, 2019]
        \textit{2019 Refinement to the 2006 IPCC Guidelines for National
        Greenhouse Gas Inventories.} E.\ Calvo Buendia, K.\ Tanabe,
        A.\ Kranjc, J.\ Baasansuren, M.\ Fukuda, S.\ Ngarize, A.\ Osako,
        Y.\ Pyrozhenko, P.\ Shermanau, S.\ Federici (eds). Published: IPCC,
        Switzerland.\\
        \url{https://www.ipcc.ch/report/2019-refinement-to-the-2006-ipcc-guidelines-for-national-greenhouse-gas-inventories/}

    \item[Martel et al., 2018]
        M.\ Martel, A.\ Glenn, H.\ Wilson, and R.\ Kröbel.
        Simulation of actual evapotranspiration from agricultural landscapes
        in the Canadian Prairies.
        \textit{Journal of Hydrology: Regional Studies}, 15, 105--118.\\
        \url{https://doi.org/10.1016/j.ejrh.2017.11.010}

    \item[Martel et al., 2021]
        M.\ Martel, A.\ Glenn, H.\ Wilson, S.\ Danielescu, R.\ Kröbel,
        W.\ Smith, B.\ McConkey, G.\ Guest, and H.\ Janzen.
        A parsimonious water budget model for Canadian agricultural conditions.
        \textit{Journal of Hydrology: Regional Studies}, 36, August 2021,
        100846.\\
        \url{https://doi.org/10.1016/j.ejrh.2021.100846}

    \item[Turc, 1961]
        L.\ Turc.
        Water requirements assessment of irrigation, potential
        evapotranspiration: simplified and updated climatic formula.
        \textit{Ann.\ Agron.}, 12, 13--49. (In French.)

\end{description}
